\chapter{Experimental or Analytical Results}

The raw output files from LS-DYNA were extracted and processed via Python 
to map the multi-phase mechanical behaviour of the construction helmets 
under the 49 Joules EN 397 drop test.

\section{Phase 1: Thermodynamic Validation}

Prior to assessing any mechanical data, the global statistics 
(\textit{glstat}) were analysed to confirm thermodynamic stability 
across all four material models. Table~\ref{tab:phase1} summarises 
the key energy conservation metrics.

\begin{table}[H]
    \centering
    \caption{Thermodynamic validation summary from LS-DYNA global 
    statistics (\textit{glstat}) for all four helmet shell materials.}
    \label{tab:phase1}
    \begin{tabularx}{\textwidth}{L C C C C C}
        \toprule
        \textbf{Material} & \textbf{Initial Energy (J)} & \textbf{Final Energy (J)} & \textbf{Energy Drift (\%)} & \textbf{Hourglass (J)} & \textbf{Valid?} \\
        \midrule
        Virgin HDPE       & 49.50 & 49.49 & $-0.0218$ & 0.00 & Yes \\
        ABS Plastic       & 49.61 & 49.59 & $-0.0288$ & 0.00 & Yes \\
        Bamboo-HDPE 30\%  & 49.57 & 49.56 & $-0.0329$ & 0.00 & Yes \\
        Recycled PC (rPC) & 49.78 & 49.77 & $-0.0215$ & 0.00 & Yes \\
        \bottomrule
    \end{tabularx}
\end{table}

Across the simulation time, total energy was conserved in all four models, with a maximum 
drift of only $-0.0329$\% observed in the Bamboo-HDPE model. The 
hourglass energy remained at 0.00 J throughout all simulations.

\section{Phase 2: Macroscopic Force Transmission (EN 397)}
The core requirement of the EN 397 standard is that the transmitted force 
must remain below 5.0 kN. The raw nodal acceleration data exhibited 
high-frequency oscillations characteristic of explicit penalty contact 
formulations, as shown in Figure~\ref{fig:phase2_raw}. For example, the 
raw Virgin HDPE trace reached a peak of 8.09 kN, and the Bamboo-HDPE 
composite reached 9.91 kN.

Upon applying the 4-pole phaseless SAE J211 CFC 60 Butterworth low-pass 
filter (100 Hz cutoff), the high-frequency components were suppressed, 
as shown in Figure~\ref{fig:phase2_filtered}. The raw data 
(Figure~\ref{fig:phase2_raw}) contains oscillations exceeding 10 kN, 
while the filtered curves (Figure~\ref{fig:phase2_filtered}) show 
smooth force--time histories.

\begin{figure}[H]
    \centering
    \begin{subfigure}[b]{0.85\textwidth}
        \includegraphics[width=\textwidth]{Phase2_ForceTransmission.pdf}
        \caption{Raw unfiltered transmitted force. High-frequency 
        spikes (up to $\sim$10 kN) are numerical contact artefacts, not 
        physical failures.}
        \label{fig:phase2_raw}
    \end{subfigure}
    \vspace{1em}
    \begin{subfigure}[b]{0.85\textwidth}
        \includegraphics[width=\textwidth]{Phase2_ForceTransmission_Filtered.pdf}
        \caption{Filtered transmitted force after SAE J211 CFC 60 
        processing. The 5.0 kN EN 397 threshold is shown as a dashed 
        line.}
        \label{fig:phase2_filtered}
    \end{subfigure}
    \caption{Transmitted Force vs.\ Time for all four helmet shell 
    materials under the EN 397 49 J impact: (a) raw explicit solver 
    output and (b) after application of the SAE J211 CFC 60 low-pass 
    filter. The comparison demonstrates that signal processing is 
    required to recover physically meaningful force histories from 
    explicit penalty contact simulations.}
    \label{fig:phase2}
\end{figure}

\begin{table}[H]
    \centering
    \caption{Summary of EN 397 force transmission results for all four 
    helmet shell materials.}
    \label{tab:results}
    \begin{tabularx}{\textwidth}{L C C C C}
      \toprule
      \textbf{Material} & \textbf{Raw peak (kN)} & \textbf{Filtered peak (kN)} & \textbf{EN 397 limit (kN)} & \textbf{Pass/Fail} \\
      \midrule
      Virgin HDPE       & 8.09 & 4.05 & 5.0 & PASS \\
      ABS Plastic       & 2.74 & 2.25 & 5.0 & PASS \\
      Bamboo-HDPE 30\%  & 9.91 & 4.51 & 5.0 & PASS \\
      Recycled PC (rPC) & 8.93 & 4.54 & 5.0 & PASS \\
      \bottomrule
    \end{tabularx}
\end{table}

\subsection{Benchmarking against Existing Data}
The filtered peak force values were compared to results from existing 
studies and experimental reports. Long et al. \cite{long2015} reported 
transmitted forces of 3.2--3.8 kN for HDPE shells, while Tham et al. 
\cite{tham2014} reported experimental ranges of 3.0--4.5 kN for Type I 
HDPE shells under 49 J impacts. The current Virgin HDPE baseline (4.05 kN) 
is consistent with these recorded ranges.

The Bamboo-HDPE 30\% wt result (4.51 kN) and the Recycled PC shell (4.54 kN) 
also fall within the compliant range defined by the EN 397 standard.

\section{Phase 3: Energy Absorption Kinematics}

Shock absorption is fundamentally the physical conversion of 
kinetic energy into internal strain energy. Table~\ref{tab:phase3} 
presents the final energy partition for each material at 
$t = 0.02$ s, extracted from the \textit{glstat} output files.

\begin{table}[H]
    \centering
    \caption{Energy partition at end of simulation ($t = 0.02$ s) 
    for all four helmet shell materials. Initial impact energy 
    $\approx$ 49.5 J across all models.}
    \label{tab:phase3}
    \begin{tabularx}{\textwidth}{L C C C C}
        \toprule
        \textbf{Material} & \textbf{Internal Energy (J)} & \textbf{Residual Kinetic (J)} & \textbf{Contact Friction (J)} & \textbf{Absorption Ratio (\%)} \\
        \midrule
        Virgin HDPE       & 27.35 & 15.86 & 6.29 & 55.3 \\
        ABS Plastic       & 26.21 & 17.59 & 5.79 & 52.8 \\
        Bamboo-HDPE 30\%  & 22.84 & 20.25 & 6.47 & 46.1 \\
        Recycled PC (rPC) & 16.67 & 26.90 & 6.20 & 33.5 \\
        \bottomrule
    \end{tabularx}
\end{table}

Virgin HDPE recorded the highest proportion of impact energy 
(27.35 J, 55.3\%), with the lowest residual kinetic energy (15.86 J). 
ABS recorded 26.21 J of internal energy and 5.79 J of contact friction 
energy. 

Recycled PC recorded 16.67 J of internal strain energy and 26.90 J of 
residual kinetic energy. The contact friction energy ranged from 
5.79 J to 6.47 J across all material models. 

The mechanical response to this energy partitioning is visible in the 
resulting stress contours. Figure~\ref{fig:stress_contours} compares 
the peak Von Mises stress distribution across the candidate materials, 
illustrating the distinct load distribution mechanisms between 
ductile and high-modulus shells.

\begin{figure}[H]
    \centering
    \begin{subfigure}[b]{0.48\textwidth}
        \includegraphics[width=\textwidth]{stress_top_hdpe.png}
        \caption{Virgin HDPE}
        \label{fig:stress_hdpe}
    \end{subfigure}
    \hfill
    \begin{subfigure}[b]{0.48\textwidth}
        \includegraphics[width=\textwidth]{stress_top_rpc.png}
        \caption{Recycled PC}
        \label{fig:stress_rpc}
    \end{subfigure}
    \vskip\baselineskip
    \begin{subfigure}[b]{0.48\textwidth}
        \includegraphics[width=\textwidth]{stress_top_abs.png}
        \caption{ABS Plastic}
        \label{fig:stress_abs}
    \end{subfigure}
    \hfill
    \begin{subfigure}[b]{0.48\textwidth}
        \includegraphics[width=\textwidth]{stress_top_bamboo.png}
        \caption{Bamboo-HDPE 30\%}
        \label{fig:stress_bamboo}
    \end{subfigure}
    \caption{Peak Von Mises stress contours during top impact (EN 397). High-modulus shells (ABS, rPC) use higher flexural stiffness to distribute impact energy across a wider area of the EPS liner, while ductile materials (HDPE, Bamboo-HDPE) exhibit more localized stress concentrations at the impact site.}
    \label{fig:stress_contours}
\end{figure}

The distribution of internal strain within the shell structure further 
validates the energy absorption hierarchy. Figure~\ref{fig:strain_contours} 
contrasts the localized plastic strain in the HDPE baseline with the 
dispersed elastic strain in the rPC alternative.

\begin{figure}[H]
    \centering
    \begin{subfigure}[b]{0.48\textwidth}
        \includegraphics[width=\textwidth]{strain_section_hdpe.png}
        \caption{Virgin HDPE (Localized strain)}
        \label{fig:strain_hdpe}
    \end{subfigure}
    \hfill
    \begin{subfigure}[b]{0.48\textwidth}
        \includegraphics[width=\textwidth]{strain_section_rpc.png}
        \caption{Recycled PC (Dispersed strain)}
        \label{fig:strain_rpc}
    \end{subfigure}
    \caption{Cross-sectional internal strain distribution at peak impact. The ductile HDPE shell undergoes significant localized deformation, whereas the high-modulus rPC distributes strain across the curvature of the dome.}
    \label{fig:strain_contours}
\end{figure}

\section{Phase 4: Kinematic Displacement and Structural Crush Depth}
The final mechanical metric evaluated was the relative kinematic
displacement between the headform center of gravity and the helmet apex.
This represents the true structural crush depth, independent of the
global movement of the assembly. The corresponding peak normal stress 
(the primary compressive stress acting perpendicular to the crown impact plane along the vertical Z-axis) 
was also tracked to correlate shell stiffness with internal pressure. Results are summarized in
Table~\ref{tab:crush_depth} and Figure~\ref{fig:phase4_kinematics}.

\begin{figure}[H]
    \centering
    \begin{subfigure}[b]{0.48\textwidth}
        \includegraphics[width=\textwidth]{Phase5_CrushDepth.pdf}
        \caption{Crush Depth vs. Time}
        \label{fig:crush_depth_plot}
    \end{subfigure}
    \hfill
    \begin{subfigure}[b]{0.48\textwidth}
        \includegraphics[width=\textwidth]{Phase4_NormalStress.pdf}
        \caption{Peak Normal Stress (Vertical Axis)}
        \label{fig:normal_stress_plot}
    \end{subfigure}
    \caption{Kinematic and structural response during Phase 4: (a) relative kinematic displacement (crush depth) and (b) peak normal stress histories (compressive stress acting along the vertical Z-axis), illustrating the distinct stiffness responses between materials.}
    \label{fig:phase4_kinematics}
\end{figure}
\begin{table}[H]
    \centering
    \caption{Maximum structural crush depth (relative displacement) for 
    all four helmet shell materials under 49 J impact.}
    \label{tab:crush_depth}
    \begin{tabularx}{\textwidth}{L C}
        \toprule
        \textbf{Material} & \textbf{Max Crush Depth (mm)} \\
        \midrule
        Virgin HDPE       & 18.42 \\
        ABS Plastic       & 12.15 \\
        Bamboo-HDPE 30\%  & 15.68 \\
        Recycled PC (rPC) & 10.23 \\
        \bottomrule
    \end{tabularx}
\end{table}

The kinematic displacement data directly supports the energy partitioning mechanisms identified in Phase 3. The highly ductile Virgin HDPE model exhibited the deepest structural crush (18.42~mm), confirming its reliance on massive plastic deformation to absorb the impact energy. The Bamboo-HDPE composite followed a similar ductile trend with 15.68~mm of crush depth. Conversely, the high-modulus materials exhibited significantly less deformation. The Recycled PC model recorded the minimum crush depth (10.23~mm), a 44.5\% reduction compared to the baseline HDPE. This restricted displacement is a direct consequence of the material's high flexural stiffness, which rapidly distributes the load to the underlying EPS liner rather than absorbing it through local shell yielding.

\section{Phase 5: 30-Degree Oblique Impact}
To evaluate the rotational kinematics critical for modern head injury mitigation, a 30-degree oblique impact simulation was conducted. The peak linear and angular accelerations recorded at the headform center of gravity, alongside the peak contact pressure, are summarized in Table~\ref{tab:oblique_results}.

\begin{table}[H]
    \centering
    \caption{Peak kinematic metrics recorded at the headform center of gravity during the 30-degree oblique impact simulations.}
    \label{tab:oblique_results}
    \begin{tabularx}{\textwidth}{L C C C}
        \toprule
        \textbf{Material} & \textbf{Peak Linear (g)} & \textbf{Peak Angular (rad/s$^2$)} & \textbf{Peak Contact (kg/mm$^2$)} \\
        \midrule
        Virgin HDPE       & 87.6  & 3482 & 2.63 \\
        ABS Plastic       & 190.8 & 5122 & 2.22 \\
        Bamboo-HDPE 30\%  & 286.9 & 12044 & 2.49 \\
        Recycled PC (rPC) & 91.3  & 896 & 2.69 \\
        \bottomrule
    \end{tabularx}
\end{table}

The distinct kinematic responses in the oblique phase are driven by 
localized surface deformation at the impact site. Figure~\ref{fig:oblique_stress} 
shows the peak Von Mises stress distribution for the ABS model during 
the 30-degree oblique impact, highlighting the primary contact mechanics.

\begin{figure}[H]
    \centering
    \includegraphics[width=0.85\textwidth]{abs_oblique_stress_final.png}
    \caption{Peak Von Mises stress distribution for the ABS Plastic shell during 30-degree oblique impact, demonstrating the localized stress concentration at the shell-anvil interface.}
    \label{fig:oblique_stress}
\end{figure}
